\begin{acknowledgments}
First and foremost, I must extend my deepest, most transfinite gratitude to my primary advisor, Dr. Georg Cantor. When the sheer cardinality of my despair approached the uncountably infinite, he was always there to remind me that my struggles were merely a minor subset of a much larger, well-ordered universe of chaos. Without his guidance, this thesis would have surely collapsed into an inconsistent mess of diagonal arguments.

To Dr. Albert Einstein, I owe a massive debt of gratitude—not only for his invaluable insights into the spatiotemporal curvature of my procrastination, but for constantly reminding me that everything is relative, including my defense deadline. His profound wisdom ensured that even when my math was entirely wrong, it was at least aesthetically pleasing in four dimensions.

I am also profoundly indebted to Dr. Alfred Tarski, who patiently sat through countless hours of my incoherent rambling. He taught me the true meaning of truth, even if we ultimately concluded that it is mathematically impossible to formally define whether this thesis makes any sense within its own language.

Finally, I want to thank the local coffee shop, the Banach-Tarski bakery (which somehow managed to turn one croissant into two every morning), and the concept of hyperbolic geometry for keeping me thoroughly ungrounded throughout this entire process.

\end{acknowledgments}
